\cotefile{157-2}
\scanpage{383}\mspage{122}

\[
H_T(Y)_{\mathrm{Cat\,hot}}=\mathrm{Hom}^{\unsure{\deleted{\circ}|\circledast}}_{\mathrm{Ind\,Cat\,h.t}}(Y,H_T)
\]

Pour résumer ces lumières et perplexités sur les $C_\infty(\underline t)$

a) À \unsure{tout} $\underline t=(t_n)_{n\in\mathbb N^{\mathrm o}}\in\underline{\mathrm{Hom\,cof}}(\mathbb N^{\mathrm o},T)$ est associé un système inductif, indexé par $\mathbb N$, de catégories $C_n(\underline t)$ $(=C_{t_n})$, et leur limite inductive

\[
C_\infty(\underline t).
\]

b) Ce\deleted{\unsure{tte construction}}\inserted{système} inductif, et aussi sa limite, est \inserted{foncteur} contravariant en $\underline t$. Si $\underline t,\underline t'\in\mathrm{Ob}\,\underline{\mathrm{Hom\,cof}}$, alors on peut comparer $C_\infty(\underline t)$, $C_\infty(t')$ et les syst. inductifs correspondants, grâce à

\diagram{\begin{tikzcd}[ampersand replacement=\&]
\& \underline t\unsure{*|\times}\underline t' \arrow[dl]\arrow[dr] \& \\
\underline t \& \& \underline t'
\end{tikzcd}}

c) L'objet

\[
\big(C_n(\underline t)\big)_{n\in\mathbb N}\in\mathrm{Ob}\,\mathrm{Ind\,Cat\,hot}
\]

défini par le système inductif $(C_n(\underline t))_{n\in\mathbb N}$ de $\mathrm{Cat}$ est \inserted{(à isom unique près)} indépendant du choix de $\underline t$, soit $H_{\unsure{T|C}}$\inserted{, ou $H$}, Les hom

\[
\underline C_\infty(\underline t')\xrightarrow{u^*}\deleted{\unsure{C}}C_\infty(\underline t)
\tag*{$(\unsure{*|\leftrightarrow})$}
\]

associé à $u:\underline t\longrightarrow\underline t'$ sont compatibles avec cette donnée, i.e. on a comm. \deleted{\unsure{du}}

\diagram{\begin{tikzcd}[ampersand replacement=\&,column sep=small]
\underset{\mathrm{Ind\,Cat\,h.t}}{\underset{\longrightarrow}{\lim}}\;\underline C_*(\underline t') \arrow[rr] \& \& \underset{\mathrm{Ind\,Cat\,h.t}}{\underset{\longrightarrow}{\lim}}\; C_{*}(\underline t)\\
\& H\deleted{\unsure{T}}\arrow[ul,"\sim"']\arrow[ur,"\sim"] \&
\end{tikzcd}}

% archive sheet number: 382
\scanpage{384}\mspage{123}

d) Il s'ensuit que les hom (*) sont des \uline{iso} de $\mathrm{Hot}$, donc tous les $C_\infty(\underline t)$ sont isomorphes dans $\underline{\mathrm{Hot}}$. L'astuce du $t*t'$ (dans le cas où on a une loi $t*t'$ associative et admettant l'objet final comme \inserted{objet} unité) donne $\simeq$ des isom \uline{canoniques}

\[
C_\infty(t)\xrightarrow{\sim} C_\infty(t')\qquad\text{dans }\mathrm{Hot}
\]

et c'est là un système transitif d'iso. Donc on a dans $\mathrm{Hot}$ un objet \deleted{$C_\infty$, défini} canonique

\[
C_\infty\in\mathrm{Ob}\,\mathrm{Hot},
\]

dont la relation à $H\in\mathrm{Ob}\,\mathrm{Ind\,Hot}$ est déduite de c), c'est une flèche canonique

\[
H\longrightarrow C_\infty\qquad\text{dans }\mathrm{Ind\,Hot}.
\]

% slanted
\margin{\uline{NB} $H$ n'est pratiquement jamais ess. constant, donc la flèche \deleted{\unsure{préci|préc.}} ci-contre n'est pas un iso dans $\mathrm{Ind\,Hot}$.}

Cette flèche induit un iso. pour les objets correspondants de $\widehat{\mathrm{Hot}_0}$\,; \unsure{Quand} on regarde $H$ \inserted{non} comme \deleted{fonct.} objet de $\widehat{\mathrm{Hot}_0}$, mais prenant le foncteur \inserted{$H_0$} induit sur $\mathrm{Hot}_0$ par le foncteur ind-représentable $H$ sur $\mathrm{Hot}$. Ce foncteur induit est ``représentable'' par l'objet \inserted{$C_\infty$} de la catégorie plus grande

\[
H_0\simeq C_\infty\qquad\text{dans }\widehat{\mathrm{Hot}_0},
\]

(si le foncteur \inserted{can.} $\mathrm{Hot}\longrightarrow\widehat{\mathrm{Hot}_0}$ est pl. fid, alors cette formule décrit $C_\infty$ \deleted{\unsure{comme}} objet de $\mathrm{Hot}$, à isom unique près, en termes de $H$, \deleted{donc} \inserted{aussi} en termes de $H\in\mathrm{Ob}\,\mathrm{Ind\,Cat\,hot}$.

% archive sheet number: 383
\scanpage{385}\mspage{124}

e) J'ignore par contre si les flèches

\[
C_\infty(\underline t')\unsure{-|\longrightarrow}\deleted{\unsure{C}}C_\infty(t)
\]

sont des iso dans $\mathrm{Cat\,hot}$, bien qu'elles soient définies en termes d'hom. de systèmes inductifs de catégories, qui sont des iso. dans $\mathrm{Ind\,Cat\,hot}$. Il en est ainsi, même pour les flèches \deleted{\unsure{entre}}

\[
C_\infty(\underline t),C_\infty(\underline t')\rightrightarrows C_\infty(t\times t')
\]

utilisées pour comparer $C_\infty(t)$ et $C_\infty(t')$ canoniquement. \unsure{Ainsi|Aussi}, je ne vois pas comment définir des flèches en sens inverse dans $\mathrm{Cat}$, qui seraient candidates à être des inverses dans $\mathrm{Cat\,hot}$.

f) Dans le cas ``géométrique'' ($\unsure{\overline{\mathrm{di}}|\overline{\mathrm{Di}}}=\underline{\mathrm{Ty}}$, situation décrite en termes des $X,\overline\Sigma$, \deleted{\unsure{$a,b\in$}} et $a,b\in\mathrm{Ob}\,X$), d'où des $C_\infty(\underline t;a,b)$, $C_\infty(a,b)$, $H(a,b)$ etc, on n'a \uline{pas} \inserted{d'op.} de composition des chemins au niveau des \uline{catégories} $C_\infty(\underline t;a,b)$ et $C_\infty(t;b,c)$, mais seulement au niveau des objets \deleted{$C_\infty(a,b)$.} $H(a,b)$, $H(b,c)$ de $\mathrm{Ind\,Cat\,hot}$ ou de $\mathrm{Ind\,Hot}$\unsure{,|.} Suis-je seulement sûr de savoir définir \uline{dans Hot}

\[
\underline C_\infty(a,b)\times C_\infty(b,c)\longrightarrow C_\infty(a,c)\qquad ?
\]

Sûrement, il n'y a pas de \unsure{problèmes|problème} \deleted{pour} les images dans $\widehat{\mathrm{Hot}_0}$, puisque là ce sont les \unsure{$\simeq$|mêmes} que celles des $H(a,b)$, $H(b,c)$, $H(a,c)$ respectivement, et on a alors accouplements \unsure{parfaits}, associatifs, avec unité, et avec inverses \deleted{pour} là où il doit y

% archive sheet number: 384
\scanpage{386}\mspage{125}

en avoir. \deleted{\unsure{C'est}} Mais pas pour les $C_\infty(a,b)$ eux-mêmes ?

C'est un peu vexant qu'avec tout ce travail, il semblerait qu'on n'ait pas même été \unsure{foutus|foutu} de définir \inserted{l'analogue dans $\mathrm{Cat}$} \unsure{des classiques espaces|du classique espace} \unsure{des|de} lacets $\Omega(X,a)$ comme un groupe \uline{dans Hot}. (Que fait Thomason ?) Mais voilà une idée \unsure{pendable|pensable|possible}. Pour définir \uline{quand même}

\[
C_\infty(a,b)\times C_\infty(b,c)\unsure{-|\longrightarrow} C_\infty(a,c)
\]

dans $\mathrm{Hot}$, on va décrire les deux facteurs $C_\infty(a,b)$ et $C_\infty(b,c)$ par des systèmes projectifs différents $\underline t,\underline t'$. On \deleted{\unsure{prendra disons}} \inserted{peut l'expliciter} en partant d'un $\tau_0$ de longueur \unsure{utile} $\geq2$, et en prenant dans les deux cas

\[
t_n=t'_n=\tau_0^n,
\]

mais des morphismes de transition différents

\[
\begin{aligned}
(t_{n+p}\longrightarrow t_n)&=\pi_p*\mathrm{id}_{\tau_0^n}\\
(t'_{n+p}-t'_n)&=\mathrm{id}_{\tau_0^n}*\pi_p
\end{aligned}
\]

où $\pi_p:\tau_0^p-1$ est la \deleted{\unsure{can}} dégénérescence totale.

On \inserted{\unsure{donc}} définira $C_\infty(a,c)$ en termes de $\underline t''\inserted{=\underline t*\underline t'}$

\[
\left\{\begin{aligned}
t''_n&=\tau_0^{2n}\\
(t''_{n+p}-t''_n)&=\pi_p*\mathrm{id}_{\tau_0^{2n}}*\pi_{\unsure{p}}.
\end{aligned}\right.
\]

Cette fois, les \deleted{morphismes} \inserted{foncteurs} de composition

\[
\deleted{\unsure{H}}C_n(\deleted{\unsure{t_n}}\underline t;a,b)\times C_n(\underline t';\unsure{b,c})\longrightarrow C_n(\underline t'';a,c)
\]

sont compatibles avec les foncteurs transition pour $n$ variable, et on trouve

% archive sheet number: 385
\scanpage{387}\mspage{126}

\[
C_\infty(\underline t;a,b)\times\underline C_\infty(\underline t';b,c)\unsure{-|\longrightarrow}\underline C_\infty(\underline t'';a,c)
\]

donc une flèche dans $\mathrm{Hot}$

\[
\mathrm{comp}_{\tau_0}:C_\infty\times C_\infty- C_\infty
\]

Il faut voir qu'elle ne dépend pas du choix de $\tau_0$, mais si on a $\tau_0$ et $\tau'_0$, on les compare par

\diagram{\begin{tikzcd}[ampersand replacement=\&,row sep=0pt]
\& \tau_0\\
\tau_0\unsure{\times|*}\tau'_0\arrow[ur]\arrow[dr] \&\\
\&\tau'_0
\end{tikzcd}}

et il suffit donc de voir que s'il existe une flèche $\tau_0\longrightarrow\tau'_0$, alors $\mathrm{comp}_{\tau_0}=\mathrm{comp}_{\tau'_0}$. Mais une telle flèche induit \deleted{un diagramme commutatif} \inserted{des flèches} \deleted{\unsure{dans $\underline{\mathrm{Hom\,cof}}$}}, $\underline t_{\unsure{\mathrm I_0|\tau_0}}\longrightarrow\underline t_{\tau'_0}$, $\underline t'_{\tau_0}\longrightarrow\underline t'_{\tau'_0}$, $\underline t''_{\tau_0}\longrightarrow\underline t''_{\tau'_0}$, d'où un diagramme dans $\mathrm{Cat}$

\diagram{\begin{tikzcd}[ampersand replacement=\&,row sep=small,column sep=small,execute at end picture={\draw[thin] (leftupper.south west) -- (trialprime.north east); \draw[thin] (leftlower.south west) -- (rightupper.north east); \draw[thin] (leftmiddle.south west) -- (trial.north east);}]
\& |[alias=trial]|{\deleted{\unsure{t_{\tau_0}}}}\&|[alias=trialprime]|{\deleted{\unsure{t'_{\tau_0}}}}\&\\
\&|[alias=leftupper]|{\deleted{t''_{\tau_0}}}\arrow[rr,leftrightarrow]\arrow[dl]\arrow[dd]\&\&|[alias=rightupper]|{\deleted{t''_{\tau'_0}}}\arrow[dl]\arrow[dd,no head]\\
|[alias=leftmiddle]|{\deleted{\underline t_{\tau_0}}}\arrow[rr,leftrightarrow]\&\&{\deleted{\underline t_{\tau'_0}}}\&\\
\&|[alias=leftlower]|{\deleted{\underline t'_{\tau_0}}}\arrow[rr,leftrightarrow]\&\&{\deleted{t'_{\tau'_0}}}
\end{tikzcd}}

\diagram{\begin{tikzcd}[ampersand replacement=\&,column sep=small]
\deleted{\unsure{H}}C_\infty(\underline t_{\tau'_0};a,b)\times C_\infty(\underline t'_{\tau'_0};b,c)\arrow[r,"{\mathrm{comp}_{\tau'_0}}"]\arrow[d] \& C_\infty(t''_{\tau'_0};a,c)\arrow[d]\\
C_\infty(\underline t_{\tau_0};a,b)\times C_\infty(t'_{\tau_0};b,\unsure{c|a})\arrow[r,"{\mathrm{comp}_{\tau_0}}"] \& C_{\deleted{\unsure{\infty}}\infty}(\deleted{\unsure{t}}t''_{\tau_0},a,c)
\end{tikzcd}}

dont on vérifie tout de suite qu'il est commutatif. Or ce diagramme peut se réécrire dans \deleted{\unsure{Cat}} $\mathrm{Hot}$ comme

\diagram{\begin{tikzcd}[ampersand replacement=\&]
C_\infty(a,b)\times C_\infty(b,c)\arrow[r,"{\mathrm{comp}_{\tau'_0}}"]\arrow[d,"\mathrm{id}"]\&C_\infty(a,c)\arrow[d,"\mathrm{id}"]\\
C_\infty(a,b)\times C_\infty(b,c)\arrow[r,"{\mathrm{comp}_{\tau_0}}"]\&C_\infty(a,c),
\end{tikzcd}}

% archive sheet number: 386
\scanpage{388}\mspage{127}

d'où $\mathrm{comp}_{\tau_0}=\mathrm{comp}_{\tau'_0}$.

Je recule devant la vraie difficulté -- celle de prouver l'associativité de cette loi de composition. Elle devient associative dans $\widehat{\mathrm{Hot}_0}$ -- mais \deleted{\unsure{pour que deux}} \inserted{pour que deux} flèches dans $\mathrm{Hot}$ soient égales (ici deux \deleted{de} \deleted{$C_\infty\times C_\infty\times C_\infty$}

\[
C_\infty(a,b)\times C_\infty(b,c)\times C_\infty(c,d)\rightrightarrows C_\infty(a,d)),
\]

suffit-il qu'il en soit ainsi de leurs images dans $\widehat{\mathrm{Hot}_0}$ -- en d'autres termes, le foncteur

\[
\mathrm{Hot}\unsure{-|\longrightarrow}\widehat{\mathrm{Hot}_0}
\]

est-il (en plus d'être conservatif) au moins fidèle (sinon pleinement fidèle) ?

Il y a aussi que cette loi de comp. est compatible avec celle des $H(a,b)$ de $\mathrm{Ind\,Hot}$, pour les flèches canoniques

\[
H(a,b)\longrightarrow C_\infty(a,b)\qquad\text{dans }\mathrm{Ind\,Hot}.
\]

Si on savait que c'est des \unsure{épis|épi}, et le restent par produits cartésiens, on aurait gagné.

Bien sûr, \deleted{\unsure{c'est}} on saura que c'est associatif, une fois qu'on saura que \inserted{la composition définie ici} est compatible avec la composition des chemins dans le cadre $1/2$ simplicial. Mais comment fait-on dans ce cadre, pour composer les chemins de longueur \deleted{\unsure{(et de t|et de t}} non bornée ? N'y a-t-il pas là exactement la même difficulté ? À moins de \unsure{s'app|s'en}

% archive sheet number: 387
\scanpage{389}\mspage{128}

tenir en \unsure{se ramenant|se restreignant|raisonnant} tout de suite, pour la théorie des opérations $\Omega$, aux complexes de Kan, \deleted{\unsure{Mais}} et aux chemins de type $\Delta_+$. Mais pour les composer \uline{fonctoriellement}, il me semble y avoir toujours le même problème -- on trouve des chemins de type $\Delta_+\Delta_+$, \inserted{$=\Delta_2^{\unsure{+|\to}}$} et non pas $\Delta_+$, -- mais on peut alors \deleted{\unsure{contracter, pour obtenir}} \inserted{composer les homotopies} directes (revient à s'induire sur $\Delta_+\simeq\{0,2\}\subset\Delta_2^+$).

Or j'ai fait le \unsure{formalisme} des catégories de chemins sans me soucier des conditions \unsure{du|de} type Kan. Je présume que dans le cadre semi-simplicial, les $\underline{\mathrm{Hom}}(\Delta_1^{\unsure{+|\to}},X)$ ne donnent \uline{pas} le \unsure{bon} objet de chemins, pour $X$ pas de Kan. (Car sinon, il semblerait que pour $X=X_{\unsure{\bullet|0}}$ à composantes \uline{finies}, le calcul des $\pi_{\unsure{i|1}}$ devrait être fini, ce qui n'est pourtant pas le cas, m'a-t-on toujours assuré \ldots)

Je présume pourtant, quand on prend $\overline\Sigma=\mathrm{Fl}(X)$, que les \deleted{catégories} objets $\underline{\mathrm{Ch}}_\infty(X;a,b)$ et $\underline{\mathrm{Ch}}_\infty(\overline X;a,b)$ \deleted{\unsure{-}} sont \deleted{\unsure{pas}} \inserted{\deleted{en général}} isomorphes dans $\mathrm{Hot}$, quand on immerge $X$ dans une enveloppe de Thomason\unsure{-| }Kan $\overline X$. \deleted{Cela \unsure{me mènerait à|me ferait} \unsure{dire|donc} que la question d'associativité pourtant,}

% archive sheet number: 388
\scanpage{390}\mspage{129}

11. Rapide retour \deleted{\unsure{sur}} sur les $(M,\overline\Sigma)$, dans l'optique des dérivateurs.

Quand on s'intéresse au formalisme des $\underline{\mathrm{Ch}}_{\overline\Sigma}(X;a,b)$ en vue d'un \uline{dérivateur} $\mathbb D_{X,\overline\Sigma}$ qu'on voudrait \unsure{définissable} en termes de $X,\overline\Sigma$, il semblerait que le \deleted{\unsure{qu'on}} prédérivateur qu'on obtient bel et bien ne mérite le nom de dérivateur que lorsqu'on suppose que pour toute flèche

\[
Y\longrightarrow Y'
\]

dans $\mathrm{Cat}$ qui est une $\mathrm{Hot}$-équivalence, les flèches correspondantes

\[
\deleted{\unsure{\mathrm{Ch}|\mathrm{Ob}}_{\unsure{\infty}}}\mathrm{Hom}_{\mathrm{Ind\,Cat\,hot}}(Y',H(a,b))\longrightarrow\mathrm{Hom}_{\mathrm{Ind\,Cat\,hot}}(Y,H(a,b)),
\]

ou encore

\[
\underline{\lim}\,\mathrm{Hom}_{\mathrm{Cat\,hot}}(Y',\mathrm{Ch}_{n,\overline\Sigma}(a,b))\longrightarrow
\underset{n}{\underline{\lim}}\,\mathrm{Hom}_{\mathrm{Cat\,h.t}}(Y,\mathrm{Ch}_{n,\overline\Sigma}(a,b))
\]

sont des bijections. J'ai regardé \inserted{dans \unsure{VI|VII}} ce que ça signifie pour l'objet $H(a,b)\unsure{=|\simeq}(\mathrm{Ch}_{n,\overline\Sigma}(a,b))_n$ de $\mathrm{Ind\,Cat\,hot}$. C'est cette condition seulement qui assure, du point de vue des dérivateurs, une interprétation convaincante de $\mathrm{Ch}_{\infty,\overline\Sigma}(a,b)$

\[
\mathrm{Ch}_{\infty,\overline\Sigma}(a,b)=\underset{n}{\underset{\longrightarrow}{\lim}}\mathrm{Ch}_{n,\overline\Sigma}(a,b)
\]

en tant qu'objet de $\mathrm{Hot}$, ou plus exactement de $\widehat{\mathrm{Hot}_0}$, car on a \inserted{alors} pour $Y$ une catégorie \uline{finie}

\[
\underset{\substack{\Vert\\ \displaystyle\mathrm{Hom}_{X(Y),\overline\Sigma_Y^{-1}}(a_X,b_X)}}{\mathrm{Hom}^{\deleted{\unsure{\circ}}}_{\mathrm{Cat\,hot}}(Y,H(a,b))}
\simeq\mathrm{Hom}_{\mathrm{Hot}}(Y,\underline{\mathrm{Ch}}_{\infty,\overline\Sigma}(a,b)).
\]

% archive sheet number: 389
\scanpage{391}\mspage{\deleted{\unsure{129}}\unsure{130|129}}

\deleted{\unsure{ne se résoudra sans doute}}

Car je veux quand même que le type dans $\mathrm{Hot}$ soit le bon, qu'il corresponde bien à ``l'homotopie'' des espaces de chemins. Donc \deleted{que} on \unsure{voudra|voudrait} qu'un morphisme

\[
X\xrightarrow{f}X'
\]

dans $(\deleted{\unsure{\mathrm{Hot}}}\mathrm{Cat})$, qui est une $\mathrm{Hot}$-équivalence, définisse des

\[
\underline{\mathrm{Ch}}_\infty(X\unsure{,|;}a,b)\unsure{-|\longrightarrow}\underline{\mathrm{Ch}}_\infty(X',f(a),f(\unsure{a|b}))
\]

qui soient eux aussi des $\mathrm{Hot}$-équivalences.

D'autre part, c'est évidemment compatible avec la \unsure{brillante} composition des chemins, que j'ai réussi à définir \unsure{entre} ces objets de $\mathrm{Hot}$. Donc \deleted{pour} (sous réserve d'avoir vérifié ce que je viens de dire) la vérification de l'associativité se ramènerait au cas où $X$ est de Thomason-Kan. Mais dans ce cas, il y a des chances que, tout comme dans le cas $1/2$ simplicial, les $\underline{\mathrm{Ch}}_\infty(X;a,b)$ peuvent se définir \deleted{\unsure{en ter}} (comme objets de $\mathrm{Hot}$) comme étant simplement les chemins de type $\Delta_+$ de $a$ \unsure{à} $b$, qu'on compose alors de \unsure{façon} triviale. Mais non, ces catégories sont discrètes !

Je renonce pour l'instant à poursuivre \deleted{\unsure{plus}} la question d'associativité. Demander à Thomason pour les propriétés du foncteur $\mathrm{Hot}\longrightarrow\widehat{\mathrm{Hot}_0}$.

% archive sheet number: 390
\scanpage{392}\mspage{\deleted{\unsure{130}}\unsure{131|130}}

Cela signifie, à très peu de choses près, qu'on a une interprétation de $\underline{\mathrm{Ch}}_{\infty,\overline\Sigma}(a,b)$, en tant qu'objet de $\mathrm{Hot}$, comme un ``$\mathrm{Hom}$ externe'' \inserted{$\mathrm{Hom\,ext}(a,b)$} à valeurs dans $\mathrm{Hot}$, défini sur la catégorie \inserted{$M=$} $X$ à l'aide du seul ensemble admissible \deleted{\unsure{\Sigma}} de flèches $\overline\Sigma\subset\mathrm{Fl}(\deleted{X}M)$, \inserted{opération} satisfaisant la condition

\[
\left|\enspace\underline{\mathrm{Hom}_{\mathrm{Hot}}(Y,\deleted{\unsure{\mathrm{Ch}}}\mathrm{Hom\,ext}(a,b))\simeq\mathrm{Hom}_{\deleted{X}M(Y)\overline\Sigma_Y^{-1}}(a_X,b_X)}\enspace\right|
\tag*{(*)}
\]

(qui peut s'interpréter aussi, soit comme

\[
\mathrm{Hom}_{\deleted{X}M}(p_!(a_{\unsure{Y|X}}),b)
\]

si $p_!:\underbrace{M(Y)\overline\Sigma_Y^{-1}}_{\mathbb D(Y)}\inserted{\longrightarrow M}$ existe, soit comme

\[
\mathrm{Hom}_{M}(a_{\deleted{\unsure{X|Y}}},p_*(b_{\unsure{Y|X}}))
\]

si $p_*:\underbrace{M(Y)\overline\Sigma_Y^{-1}}_{\mathbb D Y}\longrightarrow\underset{\substack{\Vert\\ \mathbb D(e)}}{M}$ existe)\,; avec le seul petit ennui que cette formule n'est valable que pour $Y$ \uline{\unsure{fini|finie}}. Si on veut exprimer le second membre pour $Y$ variable quelconque dans $\mathrm{Cat}$, il ne suffit pas d'un objet de $\mathrm{Hot}$ ($\mathrm{Hom\,ext}(a,b)$) dans le premier membre, mais il faut le Ind-objet $H(a,b)=(\mathrm{Ch}_n(a,b))_{n\in\mathbb N}$ dans $\mathrm{Hot}$. Si donc on tient à travailler \deleted{avec des catégories} \inserted{dans un ``domaine''} \inserted{$\mathrm{Diag}$} du dérivateur $\mathbb D=\mathbb D_{M,\unsure{\Sigma|\overline\Sigma}}$ aussi grand que possible, en se bornant \deleted{\unsure{pour}} à des catégories finies

% archive sheet number: 391
\scanpage{393}\mspage{132}

et qu'on veut qu'une formule (*) reste valable \inserted{pour tt $Y$ dans le domaine}, il faut ``en toute rigueur'' définir le $\mathrm{Hom\,ext}(a,b)$ non comme objet de $\mathrm{Hot}$, mais de $\mathrm{Ind\,Hot}$, \unsure{Il} se décrirait donc alors comme

\[
\mathrm{Hom\,ext}(a,b)=(\mathrm{Ch}_{n,\overline\Sigma}(a,b))_n
\quad\left|\begin{gathered}\text{Ind-objet}\\\text{dans }\mathrm{Hot},\end{gathered}\right.
\]

et l'objet $\mathrm{Ch}_{\infty,\overline\Sigma}(a,b)$ de $\mathrm{Hot}$ en est une sorte de limite inductive.

Je me suis demandé si la condition préalable, savoir que les \inserted{foncteurs} $H(a,b)$ \deleted{\unsure{fonc}} sur $\mathrm{Cat}$, transforment $\mathrm{Hot}$-équivalences $Y'\unsure{-|\longrightarrow} Y$ en bijections, était bel et bien une (forte) restriction. Je me suis convaincu que c'est bien le cas, déjà dans le cas particulier où $\overline\Sigma=\mathrm{Fl}(X)$. Il me semble que lorsque \inserted{la comp. connexe de} $\Omega(X,a)$ n'est \deleted{\unsure{belle}} pas contractile, la dite condition n'est pratiquement jamais satisfaite. J'ai pensé \inserted{p.ex.} au cas \deleted{d'une} \inserted{\unsure{où|d'où} \deleted{$\mathrm{Cat}$}} $X$ a le type \deleted{topologique} d'homotopie (dans $\mathrm{Hot}$) d'une $2$-sphère, de sorte que son $\Omega(X,a)=\mathrm{Ch}_\infty(X;a,a)$ doit avoir un $\pi_1\simeq\mathbb Z$. Prenons $Y$ ayant le type \unsure{dans|d'hos} $\mathrm{Hot}$ d'un cercle, disons

\diagram{\begin{tikzcd}[ampersand replacement=\&,column sep=small,row sep=small]
\&\&a\&\\
Y:\&\unsure{o|0}\arrow[ur]\arrow[dr]\&\&\unsure{o'|\bullet'|0'}\arrow[ul]\arrow[dl]\\
\&\&b\&\phantom{o},
\end{tikzcd}}

% archive sheet number: 392
\scanpage{394}\mspage{133}

donc les foncteurs

\[
\unsure{\lambda|f}:Y\longrightarrow\mathrm{Ch}_\infty(X,a,a)\overset{\mathrm{def}}{=}\Omega(X,a)
\]

correspondent à \uline{certains} \deleted{chemins} \inserted{``lacets''} de $\Omega(X,a)$ (mais lacets à origine non précisée), ils ont un type d'homotopie \inserted{\deleted{\unsure{au sens}} topologique} (``nombre de tours'' autour de la sphère $X$) donné par un entier $n\in\mathbb Z$ \inserted{(le ``degré'' du lacet $\lambda$)}. L'intuition évidente, c'est que l'ensemble des degrés \inserted{$\delta$} qu'on peut obtenir ainsi, avec une catégorie \uline{finie} donnée $Y$ (qui représente \deleted{le \inserted{esp} cercle \unsure{topo}} le cercle $S^1\in(\mathrm{Hot})$), est fini -- \inserted{\unsure{$\lvert\delta\rvert$|!!}} majoré \deleted{\unsure{d'homotopie} $\mathrm{Cat}$ \deleted{\unsure{donc des}}} sûrement par le nombre des \deleted{\unsure{ses|ces}} flèches (donc \unsure{au moins|au mieux} par $4$ dans le cas envisagé, chaque flèche \inserted{de $Y$} faisant faire \uline{au plus} un tour autour de \deleted{\unsure{S}} du pseudo-cercle $\Omega(S^2,a)$).

% slanted
\margin{probablement on \unsure{ne|en} trouve tout au plus que les degrés $\delta$ entre $-2$ et $+2$.}

D'autre part, quand on représente $S^1$ par des catégories $Y_n^{\deleted{\unsure{*}}}$ \deleted{\unsure{multiplicité}} arbitrairement grandes (disons déduites des \deleted{\unsure{\Omega}} intervalles $\unsure{\mathscr L|\mathcal L}_n=(\Delta_+\Delta_-)^n$ en identifiant \deleted{l'origine} \inserted{la} source et le but), on doit trouver des degrés arbitrairement élevés. Or toute $Y_n$ s'envoie dans $Y$ $(=Y_2)$ par une \deleted{\unsure{ff}} \inserted{flèche}

\[
Y_n=Y'\unsure{-|\longrightarrow}Y=Y_2
\]

% archive sheet number: 393
\scanpage{395}\mspage{134}

qui est de degré $1$, donc qui est une $\mathrm{Hot}$ équivalence. \deleted{Donc} \inserted{Mais} l'application correspondante

\[
\pi_0\underline{\mathrm{Hom}}(Y,\Omega(X,a))\unsure{-|\longrightarrow}\pi_0\underline{\mathrm{Hom}}(Y',\Omega(X,a))
\]

n'est pas bijective, Ce genre d'argument me semble devoir marcher chaque fois que l'un des groupes d'homotopie \deleted{\unsure{du}} $\pi_j(X,x)$, pour $j\geq2$ (donc l'un des groupes d'homotopie $\pi_i$ de $\Omega(X,x)$, pour $i\geq1$) est \uline{infini}, en prenant pour $Y,Y'$ \deleted{\unsure{deux}} différents représentants \inserted{$Y_n$} de la sphère $S^i$, \deleted{\unsure{suivant}} disons \deleted{de la forme $\unsure{\mathscr L|\mathcal L}_n$} les $Y_n$ déduits des \deleted{\unsure{la qu}} $\mathscr L_n$ \inserted{$(n\geq2)$} en prenant $\mathscr L_n^i$ \inserted{$=\mathscr L_n\times\mathscr L_n\,\ldots\times\mathscr L_n$} (qui est une sorte de $i$-cube combinatoire), et en identifiant son ``bord'' à un seul point. Je présume que le nombre d'éléments \inserted{de $\pi_i(\Omega(X,x))$} qu'on peut obtenir à partir de \deleted{l'application} flèches $Y_n\longrightarrow\Omega(X,x)$ est limité par la taille de $Y_n$, mais on peut \inserted{sûrement} les obtenir chacun, en prenant $n$ assez grand. Si donc $\unsure{\zeta|\xi}\in\pi_i(\Omega)$ échappe à $Y_n$ $(n_0=\unsure{n|0}$, disons), on peut \unsure{un|en} $n'$ assez grand pour que \deleted{\unsure{il}} $\zeta$ n'échappe pas à $Y_{n'}$, et \unsure{on|en} choisit

\[
Y_{n'}\longrightarrow Y_n
\]

% archive sheet number: 394
\scanpage{396}\mspage{135}

qui soit une $\mathrm{Hot}$-équivalence (p.ex. à l'aide d'une rétraction de $\unsure{\mathscr L|\mathcal L}_{n'}$ sur $\mathscr L_n$, \deleted{\unsure{$)_{\!2}$}} compatible avec sources et buts).

Ceci attire fortement l'attention sur le fait que si $Y,C$ \deleted{\unsure{deux}} deux catégories, la catégorie $\underline{\mathrm{Hom}}(Y,C)$ \uline{n'est pas} un représentant de \deleted{la catégorie} $\underline{\mathrm{Hom}}_{\mathrm{Hot}}(\mathrm{hot}(Y),\mathrm{hot}(C))$, on a seulement une \deleted{\unsure{appli}} flèche canonique

\[
\mathrm{hot}(\underline{\mathrm{Hom}}(Y,C))\longrightarrow \underline{\mathrm{Hom}}_{\mathrm{Hot}}(\mathrm{hot}(Y),\mathrm{hot}(C))
\tag*{(*)}
\]

i.e.

\[
\underbrace{\unsure{\mathrm{hom}|\mathrm{hot}}(\underline{\mathrm{Hom}}(Y,C))\times\mathrm{hot}(Y)}_{\substack{\Vert\\ \mathrm{hot}(\underline{\mathrm{Hom}}(Y,C)\times Y)}}\longrightarrow\mathrm{hot}(C)
\]

laquelle flèche est simplement déduite de

\[
\underline{\mathrm{Hom}}(Y,C)\times Y\longrightarrow C.
\]

Mais la flèche (*) n'a aucune envie d'être iso. C'est parce qu'elle ne l'est \inserted{déjà} pas dans le cas particulier $Y=\Delta^1$, essentiellement, qui nous oblige, pour faire une théorie des chemins \inserted{dans \deleted{\unsure{le}}\unsure{les} catégories} qui soit topologiquement raisonnable, de ne pas nous borner à $\Delta^1$ comme types de chemins, mais à prendre des intervalles \uline{composés} $\tau\unsure{\in|+}\unsure{\mathrm{Ty}|\underline{\mathrm{Ty}}}$.

% archive sheet number: 395
\scanpage{397}\mspage{136}

Par contre, (*) est iso quand $C$ est une catégorie de Thomason-Kan. Il suffit de voir que pour \deleted{tout} $Z$ de $(\mathrm{Cat})$, le $\mathrm{Hom}$ de $\mathrm{hot}(Z)$ dans les deux membres est \unsure{isomorphé} par la flèche déduite de (*). Or si $C$ est T.K., de même $\underline{\mathrm{Hom}}(Y,C)$ (immédiat sur la définition \inserted{de T.K.} en termes de prolongements de \deleted{\unsure{Y}}

\diagram{\begin{tikzcd}[ampersand replacement=\&,row sep=small,column sep=small]
A_0\arrow[dr]\&\& A\arrow[ll,"i"']\arrow[dl,dashed]\&\ldots\\
\& C\&\&
\end{tikzcd}}

), donc

% slanted
\margin{un peu vif ! Il faut savoir que les applications \unsure{cofibrantes|cofibrations} \inserted{\unsure{$i$|?}} de Thomason le restent après multiplication cartésienne par $Y$ \ldots}

\[
\begin{aligned}
&\mathrm{Hom}_{\mathrm{Hot}}(\mathrm{hot}(Z),\mathrm{hot}(\underline{\mathrm{Hom}}(Y,C)))\\
&\xleftarrow{\sim}\pi_0\big(\underbrace{\underline{\mathrm{Hom}}_{\mathrm{Cat}}(Z,\underline{\mathrm{Hom}}(Y,C))}_{\underline{\mathrm{Hom}}_{\mathrm{Cat}}(Z\times Y,C)}\big)
\simeq\big(\pi_0\underline{\mathrm{Hom}}(Z\times Y,C)\big),
\end{aligned}
\]

et \deleted{$\mathrm{Hom}(\mathrm{hot}(Z),$}

\[
\begin{aligned}
&\mathrm{Hom}_{\mathrm{Hot}}\big(\mathrm{hot}(Z),\underline{\mathrm{Hom}}_{\mathrm{Hot}}(\mathrm{hot}(\unsure{Y}),\mathrm{hot}(C))\big)
\simeq\mathrm{Hom}_{\mathrm{Hot}}\big(\underbrace{\mathrm{hot}(Z)\times\mathrm{hot}(Y)}_{\mathrm{hot}(Z\times Y)},\mathrm{hot}(C)\big)\\
&\underset{\text{\inserted{car $C$ est T.K.}}}{\xleftarrow{\sim}}\pi_0\big(\underline{\mathrm{Hom}}_{\mathrm{Cat}}(Z\times Y,C)\big),
\end{aligned}
\]

Ok.

On voit donc que les catégories $\mathrm{Ch}_{\infty,\overline\Sigma}(X;a,b)$ n'ont aucune tendance à être des catégories de Thomason-Kan. Pourtant, \unsure{son|ses} type dans $\mathrm{Hot}$, doit être ``le bon'', sans avoir d'abord à remplacer $X$ par une enveloppe de Th.K. $\overline X$.

% archive sheet number: 396
